% English translation of checked_pages/esq_119_039.tex; source: cote 119, site 39, PDF page 40, typescript page 37.
% AI-translated by GPT-6 Astra (OpenAI, ChatGPT Pro mode).
\scanpage{119/39}
\mspage{37}

6. One of the most interesting foundational theorems of (tame)
topology that ought to be developed would be a theorem
of ``dévissage'' (again !) of a proper tame map
of tame spaces,\frapparatus{V-FR-039-001}

\[ f : X \longrightarrow Y , \]

through a decreasing filtration of $Y$ by closed tame
subspaces $Y^{i}$, such that over the ``open strata''
$Y^{i} \setminus Y^{i-1}$ of this filtration, $f$ induces a locally
trivial fibration (from the tame viewpoint, it goes without saying). Such
a statement should further be generalized and made more precise in various
ways, in particular by requiring the existence of an analogous
\uline{simultaneous} dévissage for $X$ and a given finite family of closed
(tame) subspaces of $X$. The very notion of a locally
trivial fibration in the tame sense can also be strengthened
considerably, taking into account the fact that the open
strata $U_{i}$ are \uline{better} than spaces with a purely local tame
structure, because they are obtained as the difference of
two tame spaces, compact when $Y$ was compact. Between the
notion of a compact tame space (which can be realized as one of the
initial ``models'' in an $\mathbb{R}^{n}$) and that of a ``locally
tame'' (locally compact) space, which follows from it in a fairly obvious
way, there is a somewhat more delicate notion of a \uline{``globally}
\uline{tame'' space} $X$, obtained as the difference $X^{\wedge}\setminus Y$ of two compact
tame spaces, it being understood that no distinction is made between
the space defined by a pair $(X^{\wedge},Y)$ and that defined by a
pair $(X'^{\wedge},Y')$ obtained from it by a tame
(necessarily proper) map

\[ g : X'^{\wedge} \longrightarrow X^{\wedge} \]

inducing a bijection $g^{-1}(X) \xrightarrow{\ \sim\ } X$, taking $Y' = g^{-1}(Y)$.
Perhaps the most interesting natural example is that in which one
starts from a separated scheme of finite type over $\mathbb{C}$ or over $\mathbb{R}$, taking

\margin{39}
